\hypertarget{kickoff-prompt-what-to-paste-into-antigravity-to-start}{%
\subsection{What to paste into Antigravity to
start}\label{kickoff-prompt-what-to-paste-into-antigravity-to-start}}

\hypertarget{kickoff-prompt-first-the-thing-to-understand-before-you-paste-anything}{%
\subsubsection{First, the thing to understand before you paste
anything}\label{kickoff-prompt-first-the-thing-to-understand-before-you-paste-anything}}

\texttt{hackstreak-spec.md}, \path|hackstreak-spec-v2.md|, and
\path|spec-v3-public-release.md| are not three products.
They\textquotesingle re one design, revised three times --- each file
says explicitly what it keeps and what it overrides from the one before
it. Do not have the agent build v1 completely, then rebuild as v2, then
rebuild as v3. That throws away work for no reason, since the end
architecture was already decided before the first line of code gets
written.

What\textquotesingle s actually sequenced is the \textbf{build phases}
described inside v2 §11 and v3 §10 (Phase 1 → 1b → 2 → 3 → 4), not the
spec file numbers. Give the agent the full picture now, so it never
makes a Phase 1 choice that conflicts with where this is headed (a bad
schema shortcut that breaks the migration story later, say) --- but
scope the actual work it does tonight to Phase 1 + 1b only, and go no
further until that\textquotesingle s living in your daily routine for
real.

\hypertarget{kickoff-prompt-step-1--tonight-phase-1--1b-only}{%
\subsubsection{Step 1 --- tonight: Phase 1 + 1b
only}\label{kickoff-prompt-step-1--tonight-phase-1--1b-only}}

Open Antigravity in a fresh project folder with these five files present
(or pasted into context): \texttt{hackstreak-spec.md},
\path|hackstreak-spec-v2.md|, \path|spec-v3-public-release.md|,
\texttt{ui-ux-spec.md}, \path|packs/home-hack-lab.toml|. Then paste
this:

\begin{center}\rule{0.5\linewidth}{0.5pt}\end{center}

\begin{quote}
Read \texttt{hackstreak-spec.md}, \path|hackstreak-spec-v2.md|, and
\path|spec-v3-public-release.md| in that order. Together they describe
one evolving design, not three separate products --- each later file
supersedes the one before it wherever they conflict, exactly as each
file states at its top. Build according to the PHASES defined inside
\path|spec-v3-public-release.md| section 10, not according to which
spec file a detail came from.

Right now, build \textbf{only Phase 1 and Phase 1b}: the standalone
engine, CLI, pack system, levels/milestones, and local desktop nudging
--- nothing from Phase 2 onward (no server, no web client;
\texttt{ui-ux-spec.md} is for later and you can ignore it for this
step). Structure the code as a clean \texttt{engine/} +
\texttt{storage/} + \texttt{cli/} package, separated so a server can
wrap the same engine later without a rewrite, per
\path|spec-v3-public-release.md| §2. Python 3.11+, standard library
only (Phase 1b\textquotesingle s notification call may use
\texttt{urllib}, still no third-party dependency). Name the CLI binary
\texttt{ascend}.

Write the tests from \texttt{hackstreak-spec.md} section 11 and
\path|hackstreak-spec-v2.md| section 12 FIRST, then implement until
they pass. Install \path|packs/home-hack-lab.toml| as the sample pack
at \path|packs/community/home-hack-lab.toml| with its placeholder
\texttt{{[}vars{]}} left as placeholders, not my real Pi IP or folder
paths.

When finished, show me: the file tree, full \texttt{pytest} output, and
a demo of \texttt{ascend pack install},
\texttt{ascend mission start}, \texttt{ascend backfill},
\texttt{ascend levels}, \texttt{ascend}, \texttt{ascend 20},
\path|ascend done ... ---complete|, \texttt{ascend stats}, and
the nudge config/test command (or a clear note on what\textquotesingle s
stubbed if desktop notifications aren\textquotesingle t testable in this
environment). Ask me before building anything beyond Phase 1b.
\end{quote}

\begin{center}\rule{0.5\linewidth}{0.5pt}\end{center}

\hypertarget{kickoff-prompt-implementation-clarifications-from-final-review}{%
\subsubsection{Implementation clarifications from final
review}\label{kickoff-prompt-implementation-clarifications-from-final-review}}

These fix real inconsistencies found by re-reading all three specs
together. They override anything in \texttt{hackstreak-spec.md} or
\path|hackstreak-spec-v2.md| that conflicts, exactly as the
files\textquotesingle{} own precedence rules say --- these are just
spelled out explicitly because they were easy to miss otherwise.

\begin{enumerate}
\def\labelenumi{\arabic{enumi}.}
\tightlist
\item
  \textbf{Naming, everywhere.} v1 and v2 say \texttt{hack},
  \texttt{\textasciitilde{}/.hackstreak/}, \texttt{HACKSTREAK\_HOME},
  and \path|hackstreak/verifiers/| throughout --- including inside
  their own embedded "prompt for the agent" sections (v1 §12, v2 §13).
  All of that is superseded. The binary is \texttt{ascend}. Data lives
  at \texttt{\textasciitilde{}/.config/ascend/} (an
  \texttt{ASCEND\_HOME} environment variable overrides this, matching
  what the Dockerfile already uses for server mode). The database file
  is \path|$ASCEND_HOME/ascend.db| --- in standalone mode
  that\textquotesingle s
  \path|~/.config/ascend/ascend.db|; in the Docker
  image, \texttt{ASCEND\_HOME=/data}, so it\textquotesingle s
  \texttt{/data/ascend.db}. The verifier plugin folder is
  \path|~/.config/ascend/verifiers/|.
\item
  \textbf{Phase 1b is desktop notification ONLY, tonight.} v2 §11b
  describes a full Telegram-from-the-Pi system (a \texttt{sync-to-pi}
  command, a \texttt{nudge\_telegram.py} file, systemd timers on the Pi)
  and calls it "Phase 1b." That is superseded by
  \path|spec-v3-public-release.md| §7, which correctly narrows Phase
  1b to the local desktop notification only and moves Telegram to Phase
  2, sent by the server once one exists. Build only the desktop
  notification path tonight --- do not implement \texttt{sync-to-pi},
  Telegram, or anything cross-machine yet, even though v2 §11b describes
  it in more detail than the thing that actually supersedes it.
\item
  \textbf{Verifier paths must resolve \texttt{\textasciitilde{}}.} The
  starter pack\textquotesingle s \texttt{{[}vars{]}} includes
  \texttt{hack\_dir = "\textasciitilde{}/hack"}, substituted directly
  into \texttt{file} and \texttt{command} verifier paths (e.g.
  \texttt{\{hack\_dir\}/report.md}). A plain string substitution does
  NOT expand \texttt{\textasciitilde{}} to the home directory ---
  \path|Path("~/hack/report.md").exists()| checks a
  literal folder named \texttt{\textasciitilde{}}, not the home
  directory. The \texttt{file} and \texttt{command} verifiers must call
  \texttt{expanduser()} (or equivalent) on any path before using it, or
  the sample pack\textquotesingle s own verifiers will fail on night
  one. Add this as an explicit test case.
\item
  \textbf{\texttt{lint} must check \texttt{pack.kind}.}
  \path|hackstreak-spec-v2.md| §8b says every pack needs
  \texttt{kind = "project"} or \texttt{kind = "proficiency"}, but
  §3\textquotesingle s lint checklist never actually lists this as
  something \texttt{lint} validates. Add it: \texttt{lint} rejects a
  pack missing \texttt{kind} or with any other value.
\item
  \textbf{\texttt{GET /api/health}} (used by the Docker healthcheck in
  \texttt{docker/Dockerfile}) should be added to the route list in
  \path|spec-v3-public-release.md| §6 when Phase 2 is built --- a
  trivial 200-OK endpoint, not used tonight, just noted so it
  isn\textquotesingle t forgotten when the server is written.
\end{enumerate}

\hypertarget{kickoff-prompt-step-2--actually-verify-it-dont-just-accept-the-demo}{%
\subsubsection{Step 2 --- actually verify it, don\textquotesingle t just
accept the
demo}\label{kickoff-prompt-step-2--actually-verify-it-dont-just-accept-the-demo}}

Before trusting any of this:

\begin{enumerate}
\def\labelenumi{\arabic{enumi}.}
\tightlist
\item
  Skim the code it produces. You don\textquotesingle t need to review
  every line, but check the streak logic against the test table in
  \texttt{hackstreak-spec.md} §11 yourself --- run \texttt{pytest -v}
  and read what passed.
\item
  Run \texttt{ascend pack lint} and \path|ascend pack check-links|
  on \texttt{home-hack-lab.toml} yourself, read every
  \texttt{done\_when}, then \texttt{ascend review} it before trusting
  it as your real roadmap.
\item
  \textbf{Use it for real, daily, for at least a week or two} --- chess
  lesson, then \texttt{ascend}, then one step, then
  \texttt{ascend done}. This is the actual test. A demo that runs once
  in a chat is not the same as a tool that survives a deadline week.
\end{enumerate}

\hypertarget{kickoff-prompt-step-3--only-once-phase-1-is-genuinely-part-of-your-routine}{%
\subsubsection{Step 3 --- only once Phase 1 is genuinely part of your
routine}\label{kickoff-prompt-step-3--only-once-phase-1-is-genuinely-part-of-your-routine}}

Come back and hand the agent \path|spec-v3-public-release.md|
(architecture, already given) plus \texttt{ui-ux-spec.md} (the web
client design) together, and tell it to build Phase 2: the server, the
HTTP API, and the web PWA, per \path|spec-v3-public-release.md|
§10\textquotesingle s Phase 2 scope. Not before. A server and a phone
app for a habit you haven\textquotesingle t actually kept yet is
building the sequel before the first movie was any good.

\hypertarget{kickoff-prompt-step-4-and-beyond}{%
\subsubsection{Step 4 and beyond}\label{kickoff-prompt-step-4-and-beyond}}

Phase 3 (plugin verifiers, community pack CI, docs, publish) and Phase 4
(Telegram, real Web Push, public extras) follow the same rule:
don\textquotesingle t start the next one until the current one is
proven, not just built.
